\subsection{Limite de séparation de l'œil (Mire de Foucault)}

L'objectif de cette manipulation est de déterminer le pouvoir séparateur $\epsilon$ de l'œil.
Nous utilisons pour cela une mire de Foucault
présentant une période $a = 4 \pm 0,1 \text{ mm}$.
La distance limite de résolution $D_m$ est mesurée pour deux observateurs
dans différentes conditions de vision.

Le pouvoir séparateur est calculé par la relation relative aux petits angles,
décrite par la formule présentée en partie théorique. \eqref{eq:pouvoir_separation}

Les mesures expérimentales et les résultats calculés
sont synthétisés dans le tableau suivant :

\begin{table}[h!]
    \centering
    \begin{tabular}{|l|l|c|c|}
        \hline
        \textbf{Observateur} & \textbf{Condition} & \textbf{Distance $D_m$ (m)} & \textbf{Pouvoir séparateur $\epsilon$ (rad)} \\
        \hline
        Observateur 1 & Sans correction & $3,66 \pm 0,05$ & $1,09 \cdot 10^{-3}$ \\
        \hline
        Observateur 2 & Sans correction & $2,43 \pm 0,05$ & $1,65 \cdot 10^{-3}$ \\
        \hline
        Observateur 2 & Avec correction & $5,32 \pm 0,05$ & $0,75 \cdot 10^{-3}$ \\
        \hline
    \end{tabular}
    \caption{Mesures du pouvoir séparateur de l'œil pour différents sujets.}
    \label{tab:mire_foucault}
\end{table}

\subsubsection{Mesure du champ latéral}

Le champ latéral représente le diamètre de la zone objet visible à travers le microscope.
Cette mesure est effectuée à l'aide du micromètre-objet placé sur la platine.
Le micromètre utilisé présente une graduation où une grande division 
est subdivisée en 5 petites graduations ($1 \text{ p.g.} = 0,01 \text{ mm}$).

Les diamètres de champ $D_{champ}$ mesurés sont :
\begin{itemize}
    \item \textbf{Objectif $\times 4$} : $D_{champ} = 4,5 \text{ mm}$ ;
    \item \textbf{Objectif $\times 10$} : $D_{champ} = 1,8 \text{ mm}$ ;
    \item \textbf{Objectif $\times 40$} : $D_{champ} = 0,45 \text{ mm}$.
\end{itemize}

\subsubsection{Cercle oculaire}

Le cercle oculaire, correspondant à l'image de la pupille d'entrée par l'oculaire,
a été localisé à quelques millimètres au-dessus de la lentille de sortie.
Son diamètre a été estimé à :
\begin{equation}
    D_{co} \approx 1,5 \text{ mm}
\end{equation}

\subsection{Détermination du grossissement commercial}

Le grossissement a été mesuré expérimentalement à l'aide d'une chambre claire
en utilisant l'objectif $\times 10$.
Nous avons fait correspondre une longueur $AB$ sur le micromètre-objet
à une longueur $A'B'$ projetée sur une règle millimétrée externe.

L'oculaire utilisé pour cette mesure possède un grossissement propre de $G_{oc} = 10$.

Les données recueillies sont les suivantes :
\begin{itemize}
    \item Taille de l'objet : $AB = 10 \text{ p.g.} = 0,1 \text{ mm}$ ;
    \item Taille de l'image projetée : $A'B' = 1 \text{ cm} = 10 \text{ mm}$.
\end{itemize}

Le grossissement commercial $G_c$ est calculé par le rapport :
\begin{equation}
    G_c = \frac{A'B'}{AB} = \frac{10}{0,1} = 100
\end{equation}

À partir de ce résultat, nous déterminons la puissance intrinsèque $P_i$ du microscope
(exprimée en dioptries $\delta$) par la relation :
\begin{equation}
    P_i = 4 \times G_c = 400 \delta
\end{equation}

\subsection{Mesure des dimensions du Pluteus}

Cette dernière partie consiste à mesurer les dimensions réelles d'une larve d'oursin (Pluteus) 
en utilisant un oculaire micrométrique préalablement étalonné.

\subsubsection{Étalonnage de l'oculaire micrométrique}

L'étalonnage a été réalisé en comparant les graduations de l'oculaire
aux divisions du micromètre-objet (objectif $\times 10$).
L'observation directe a permis d'établir la relation suivante :
\begin{itemize}
    \item Une graduation de l'oculaire correspond exactement à une graduation du micromètre-objet ;
    \item La valeur d'une graduation du micromètre-objet est de $0,01 \mathrm{~mm}$.
\end{itemize}

On en déduit la valeur d'une graduation de l'oculaire, notée $C_{\mathrm{oc}}$ :
\begin{equation}
    C_{\mathrm{oc}} = 0,01 \mathrm{~mm} = 10 \mathrm{~\mu m/div}
\end{equation}

\subsubsection{Dimensions de la larve Pluteus}

Après avoir substitué le micromètre-objet par la lame du Pluteus,
nous avons relevé les dimensions caractéristiques de la larve (objectif $\times 10$).
Les mesures en divisions oculaires ($\mathrm{div}$) et leur conversion en micromètres ($\mathrm{\mu m}$) 
sont présentées ci-dessous :
\begin{table}[h!]
    \centering
    \begin{tabular}{|l|c|c|c|c|}
        \hline
        \textbf{Échantillon} & \textbf{Grand axe (div)} & \textbf{Petit axe (div)} & \textbf{Dimensions ($\mathrm{\mu m}$)} & \textbf{Moyenne ($\mathrm{\mu m}$)} \\
        \hline
        Œuf n°1 & 6 & 5 & $60 \times 50$ & 55 \\
        \hline
        Œuf n°2 & 8 & 4 & $80 \times 40$ & 60 \\
        \hline
        Œuf n°3 & 4 & 4 & $40 \times 40$ & 40 \\
        \hline
    \end{tabular}
    \caption{Dimensions des larves d'oursin mesurées à l'oculaire micrométrique.}
    \label{tab:mesures_oeufs}
\end{table}